\odiseachapter{Quinta parte}{Nausícaa}

A los quince años, durante mi primera y arrebatada lectura de la \emph{Odisea}, las diosas que aparecían en ella me inspiraban reverencia ---Hera, Atenea, Calipso y Helena, que contaba como divinidad--- y Penélope era como mi madre, que también se había pasado largas temporadas esperando a mi padre, durante los años que este estaba embarcado, aunque mi madre, con su picante sentido del humor, era bastante más divertida que la reina de Ítaca, siempre seria, tejiendo y destejiendo una mortaja. En cuanto a Circe, me daba un poco de miedo: con aquella hechicera valían pocas bromas. Lo mismo te convertía en cerdo que te mandaba de excursión al Hades. Me imponía tanto, de hecho, que no me atrevía a enamorarme de ella. No aquella primera vez, quiero decir. Todo se andaría.

De quien sí me enamoré fue de la alegre, confiada y generosa princesa Nausícaa, entre otras cosas porque desde que apareció en escena le presupuse los mismos años que yo tenía. Y aunque mi especulación era del todo debida a mis ganas de encontrar a una chica de mi edad en aquel libraco lleno de señoras imponentes, lo cierto es que no andaba desencaminado. La princesa entra en escena en el canto VI, cuando Atenea, después de dejar a Ulises durmiendo, agotado, en la playa ---acaba de varar en tierra de los feacios, tras pasar veinte días a la deriva en el mar, como había dispuesto Zeus---, se va al palacio del rey Alcínoo:

\begin{verse}
hasta la milbienlabrada alcoba se entró en que dormía\\
niña en presencia y figura igual a las siemprevivas,\\
Nausícaa, que era de Alcínoo cordialtanero la hija,
\end{verse}

Como de costumbre, nuestro traductor nos regala un par de joyas: la primera es «milbienlabrada», para describirnos la alcoba de la princesa, que Homero llama \gr{πολυδαίδαλον} (polydaídalon), compuesto del famoso \gr{πολυ} (poly) ---mil--- y \gr{δαίδαλον} (daídalon) ---labrada---. La segunda es «cordialtanero», la traducción que Tobal escoge aquí para \gr{μεγαλήτορος} (megalḗtoros) o de gran corazón. También nos dice «niña en presencia y figura igual a las siemprevivas», o sea una muchachita tan hermosa como una deidad. Me detengo en estas líneas, porque lo que mi libro decía ---en la traducción de Segalá--- era:

\begin{quote}
Penetró la diosa en la estancia labrada con gran primor en que dormía una doncella parecida a las inmortales por su natural y por su hermosura: Nausícaa, hija del magnánimo Alcínoo.
\end{quote}

Pero lo que yo leí, en la exaltación de mis quince años, era la emocionante traducción de Juanma. ¡Medio siglo antes de que la escribiera! Misterios de la relatividad general, capaz de conectar al casi niño de entonces con el casi anciano de ahora, en solo tres versos con sabor a mar.

Como era de esperar, Atenea no se ha colado en el palacio del rey solo para mirar. Le falta tiempo para empezar con sus enredos:

\begin{verse}
Se deslizó como soplo de viento al colchón de la chica,\\
y se inclinó a su cabeza, y ya con palabras le iba,\\
trasfigurándose en la hija del navicélebre Dimas\\
que de su misma edad era y de su corazón alegría;\\
a ella gemela, Atenea ojiglauca esto ya le decía:

Nausícaa, ¿cómo tu madre te trajo a ti así, tan dejada?\\
Ruedan por ahí descuidadas tus ropas lustrigalanas\\
y la tu boda está cerca, y hermosas tendrás que llevarlas,\\
y presentarlas hermosas a los que el cortejo te hagan.\\
Que de la gente esta gala bien sabes que en lenguas se anda, y\\
de ella se alegran el padre y la madre biensoberana.\\
Venga, que hay que ir a lavarlas apenas luz traiga el alba\\
---yo a la faena contigo---, para que bien sin tardanza\\
te me engalanes, que poco es el tiempo que habrás de muchacha;
\end{verse}

Todavía muchacha, como nos asegura la diosa, así que es difícil que pase de los quince. Ya le rondan los pretendientes, eso sí, su boda está cercana y yo, tumbado en una hamaca del Mar Menor, a la hora de la siesta, cuando lo único que se oía era el canto de los grillos, me moría de celos.

Princesa y todo, la diosa la manda a lavar sus ropas «lustrigalanas» (¡toma ya!) apenas amanezca. Plantado su enredo, Atenea se vuelve al Olimpo, que Homero nos presenta como un lugar luminoso y un poco aburrido:

\begin{verse}
Esto diciendo, partió Atenea ojiglauca de vuelta\\
hacia el Olimpo, que asiento es de dioses que nunca ---eso cuentan---\\
se tambalea: ni agítalo el viento ni lo ahogan tormentas\\
ni se le acercan las nieves, que limpio de nubes sus puertas\\
abre allí el cielo y por él la luz blanca del sol se pasea.\\
Allí los dioses felices sus días todos se alegran;\\
allí subió la ojiglauca después de que habló a la doncella.
\end{verse}

Nausícaa, entre tanto, se despierta bien de mañana y se va directa a su padre, para pedirle una carreta con la que ir a lavar al río, tal como le ha sugerido la divinal metomentodo en sueños:

\begin{verse}
«Oye, papá, ¿me armarías a mí una carreta ahora mismo\\
alta y rodada, que quiero llevar mis bonitos vestidos\\
a la corriente del río a lavar, que los tengo perdidos?\\
Que, igual que a mí, a ti te gusta, si estás en la junta reunido\\
con tus primeros, echarte a los cueros los paños bien limpios.\\
Y en el palacio aquí tienes contigo también cinco hijos,\\
dos ya casados, y tres en la flor de la edad aún mocitos,\\
y ellos adoran ponerse recién lavados sus linos\\
siempre que salen al baile; y de eso soy yo la que cuido».
\end{verse}

Párate un segundo, dulce lectora, suave lector, para disfrutar de estas líneas. Para empezar, Nausícaa usa la palabra \gr{πάππα} (páppa), que Tobal traduce casi literalmente por el «papá» castellano. Es todavía demasiado niña para llamar «padre» a Alcínoo. No es el caso de Telémaco, que, cuando por fin reconoce a Ulises, le llama \gr{ὦ πάτερ} (ô páter) y \gr{πάτερ φίλε} (páter phíle), «padre querido». Al bueno de Nolan lo han puesto de vuelta y media debido al «dad» que suelta Telémaco al hablar de su padre. Son ganas de fijarse en nimiedades, pero lo cierto es que en el pueblo en que crecí ---Alba no es su nombre, pero así me gusta llamarlo--- los hijos todavía trataban a sus padres de usted y empleaban, en catalán, el vocativo formal «pare». En toda la \emph{Odisea}, el término «papá» solo aparece una vez y no es por casualidad. Homero lo usa para caracterizar a la princesita zalamera pidiendo las llaves del coche.

Eso sí, no nombra los lustrigalanos atavíos que piensa lavar, sino que se refiere a las ropas de su padre y hermanos. El Poeta nos deja muy claro por qué:

\begin{verse}
Dijo ella así por pudor de nombrar esa boda lozana\\
ante su padre. Y él todo lo vio, y le volvió estas palabras:
\end{verse}

Y yo leía y no sabía cómo teletransportarme a Corfú ---había leído que el país de los feacios coincidía con la bella isla griega--- y defender a mi amada de todos aquellos moscones, innombrados, que la pretendían. El rey Alcínoo, por su parte, no es capaz de negarle nada a su pequeña:

\begin{verse}
«No, ni las mulas, mi niña, ni nada que pidas te niego.\\
Ve allí, que ya la carreta va a armártela el palafrenero\\
alta y de buena rodada, y canasta pondrá y sobrecielo».
\end{verse}

Todo el canto VI se caracteriza por una paz sutil, en contraste con la sucesión de catástrofes que han castigado al héroe ---y con la violencia que aún queda por llegar--- y por unos personajes a cuál más gentil. La escena entre Nausícaa y Alcínoo, la hemos vivido todos los que hemos tenido la suerte de tener una hija. Vuelvo a leer las líneas que me hacían vibrar de pasión juvenil hace cinco décadas y esta vez lo que se apodera de mi corazón es la nostalgia de aquellas mañanas de domingo, cuando Irene me colmaba de carantoñas antes de exponer su astuto plan, sabiendo que nada de lo que me pidiera le negaría. Unos pocos versos capaces de conectar a dos adolescentes y dos padres, separados por tres mil años. Misterios de la relatividad general. Misterios de la poesía. Einstein y Homero.

En fin, Nausícaa se sale con la suya, arma la carreta, sube a las sirvientas y se van al río. Allí se ponen a lavar con encomiable energía. Estas casi niñas convierten incluso la faena en un juego:

\begin{verse}
y junto al río remolinhondo las sueltan al pasto\\
dulce de grama; después en sus brazos del carro tomaron\\
la ropa, echándola al agua negra, y allí la pisaron\\
sobre las pilas, porfiando por ver quién pisaba más rápido.
\end{verse}

Terminados la lavandería y el almuerzo, siguen con brío y, en ausencia de pantallas, se ponen a jugar a la pelota:

\begin{verse}
Luego de que del condumio princesa y esclavas gozaron,\\
a la pelota les vino jugar, sus mantillas a un lado,\\
mientras abría canción Nausícaa candidosbrazos.\\
Y como reinadeflechas se anda Ártemis montes abajo,\\
por los alcores del alto Taigeto o por el Erimanto,\\
regocijándose entre jabalíes y ciervos belraudos,\\
y allí con ella las ninfas, las hijas de Zeus cabrilámpigo,\\
agrestes juegan, y Leto al mirarlas lo goza en su ánimo,\\
pues sobre todas se yergue su testa, su frente, y es claro\\
quién es de todas ---y hermosas son todas---, así también tanto\\
se distinguía de sus esclavas la virgen sin amo.
\end{verse}

Compara el Poeta a Nausícaa con Ártemis, no dejando duda alguna ni de su candor, ni de su virginidad, ni de su belleza, pero adelantándonos un dato importante. Ártemis, reina de la caza y los animales salvajes, no es precisamente una adolescente mustia y Homero nos va a mostrar enseguida que Nausícaa tampoco.

Huelga decir que el encuentro entre la princesa y el náufrago que ha celestineado Atenea está a punto de producirse.

\begin{verse}
Mas cuando se hizo la hora de echarse a casa de nuevo, y\\
mulas había que uncir y doblar los vestidos rebellos,\\
vino a la diosa, Atenea ojiglauca, un mejor pensamiento:\\
que despertara Odiseo y viera a la niña ojosbuenos\\
que a la ciudad de los hombres feacios vendría a traerlo.\\
Y la princesa lanzó la pelota a una criada y, al vuelo,\\
en la muchacha no dio, y de un hondo cabozo dio en medio.\\
Ellas chillaron muy largo y salió de su sueño Odiseo,\\
que, incorporándose, en mientes y pecho así está revolviendo:
\end{verse}

En lo que lleva de viaje, a Ulises casi lo descabezan los lestrigones, lo convierte Circe en cerdo, o se lo comen Polifemo y Escila. Ulises no las tiene todas consigo:

\begin{verse}
«¡Ay de mí! ¿En tierra de quién, de qué hombres habré yo caído?\\
¿Serán soberbios, brutales, ---no sé--- de lo justo no amigos,\\
o acogedores y bien temerosos del cielo en su espíritu?\\
Hay por aquí un guirigay mujeril, de muchachas, conmigo,\\
ninfas tal vez, que en los agrios picachos del monte han su nido,\\
y en las yerbosas praderas, y en los hontanares de ríos...\\
¿O es un humano linaje, aquí cerca de mí, vocivivo?\\
Ah, lo mejor es que salga y vaya a verlo yo mismo».
\end{verse}

Sale de entre las matas, se cubre con lo que puede y se dirige a las muchachas, casi en cuero picado. Las esclavas se llevan un susto tremendo ---imagínate, gentil lectora, jugando a la comba en el parque con tus amigas, cuando en esto se os echa encima un viejo fornido y desnudo, con ojos de loco--- y salen por piernas. Pero a Nausícaa, como buena Ártemis que es, le sobran valor y compasión. ¡Cómo no me iba a enamorar de ella!

\begin{verse}
Esto diciendo, el divino Odiseo salió entre las matas\\
y de la espesa arboleda arrancó a mano firme una rama\\
llena de hojas, porque sus verijas quedaran salvadas.\\
Iba cual bravo león que, en su fuerza fiado, anda y anda,\\
de lluvia herido y de viento ---sus dos ojos vivos dos llamas---,\\
y en medio de los carneros y de los bueyes se lanza\\
y de las ciervas montesas, y el vientre, incluso, le manda\\
a las ovejas en prieto redil tentar de alcanzarlas;\\
así Odiseo a juntarse con las belrubias muchachas\\
iba, aunque estaba desnudo: la necesidad lo apretaba.\\
Se les mostró aterrador, de la sal malajado y el agua,\\
y ellas huyeron acá y acullá por la playa escarpada;\\
solo la hija de Alcínoo quedó, que Atenea llenaba\\
su corazón de osadía y libraba sus miembros de alarma.\\
Y se tenía en pie frente a él; y Odiseo dudaba\\
si, a sus rodillas asiéndose, hablarle a la niña ojiclara,\\
o si de lejos mejor, con tranquilizadoras palabras,\\
pedirle por la ciudad y por si algo de ropa le daba.\\
Tal cavilando, en la duda, pensó que era cosa más sabia\\
hablar de lejos mejor, con tranquilizadoras palabras,\\
no fuera a ser que al tomar sus rodillas la joven se airara.
\end{verse}

No tiene un pelo de tonto Odiseo. Están los dos desnudos ---Nausícaa ha estado lavando sus ropas y presumiblemente los bañadores no se habían inventado todavía--- y sabe que su aspecto es terrible. No solo es viejo ---debe andar por los cincuenta---, sino que viene deplorable, malajado de sal y agua. Él mismo se siente como un mendigo, su impulso es abrazar las rodillas de la princesa y pedirle caridad. El orgulloso héroe, saqueador de ciudades, amante de deidades, reducido a la condición de pordiosero. Otra vez, Homero nos deja patidifusos, mostrándonos un ángulo del héroe que no habíamos visto hasta ahora y que nos lleva a compadecernos de él.

\begin{verse}
De hinojos, reina, ante ti, que no sé si eres diosa o doncella.\\
Que si deidad de las que hacen del ancho cielo vivienda,\\
hija de Zeus el grande, diría que tú Ártemis fueras:\\
en estatura y en formas y en trazas de estar te andas cerca.\\
Y si de humanas mortales tú eres, que habitan la tierra,\\
feliz tres veces tu padre, tres veces tu madre y tu reina,\\
y tus hermanos tres veces felices; que bien se les queman\\
sus corazones por ti de alegría, al mirar cómo entra\\
siempre a la danza una flor como tú, botón de belleza.\\
Pero de todos el más felicísimo en sus entretelas\\
aquel que a casa te lleve ganando en arras y en prendas.\\
Nunca a los ojos me vino persona mortal como esta,\\
no, ni mujer ni varón; que es mirar y el fervor me embelesa.
\end{verse}

Derrotado como está, Ulises no ha perdido la labia. O quizá sus palabras no sean simple cortesía. Puede que ---como a mi yo de entonces--- la visión de la princesa lo haya fulminado y el fervor que anuncia sea real. Pero no está en condiciones de hacer otra cosa que suplicar y esa indefensión conmueve al lector y a la princesa:

\begin{verse}
Reina, piedad ten de mí; que he llegado hasta ti la primera\\
tras de lo mucho pasado, y de nadie más sé que aquí sea\\
un morador, como tú, en esta ciudad y esta tierra.\\
Dime el poblado por dónde, y deja que me eche una tela,\\
si alguna aquí te trajiste para envolverte las prendas,
\end{verse}

Nausícaa entonces se presenta, sin pizca de miedo y hasta un poco altiva:

\begin{verse}
Y habla, mirándolo, ahora Nausícaa candidosbrazos:\\
«Puesto que a ti, forastero, ni vil se te ve ni insensato,\\
sabrás que Zeus Olímpico dicha la da a los humanos,\\
malos y buenos, según el querer que le venga a él al caso;\\
conque, si a ti eso te dio, hasta el fin deberás soportarlo.
\end{verse}

Se repite un tema ya familiar: los olímpicos en general y Zeus en particular hacen con los hombres lo que les da la gana y lo único que les queda a estos es resignarse. Enseguida le ofrece generosamente su ayuda.

\begin{verse}
Ahora, pues que a nuestra tierra y a nuestra ciudad has llegado,\\
no ha de faltarte vestido, ni nada, tampoco, de cuanto\\
ley es prestar a quien viene angustiado en pido de amparo.\\
Dónde el poblado sabrás, y el nombre de sus paisanos.\\
Esta ciudad y esta tierra asiento son hoy de feacios;\\
yo soy la hija de Alcínoo grancorazón, a su brazo\\
de los feacios el brío y poder se le ha encomendado».
\end{verse}

El tema central de la película de Nolan puede resumirse así: la ley de Zeus obliga a los hombres a la \gr{ξενία} (\emph{xenía}), esto es, a ser hospitalarios ---no por virtud, por cierto, sino por precaución, ya que el forastero al que se recibe bien podría ser un dios disfrazado---. Ulises y los aqueos toman Troya con un engaño, saltándose esa ley y precipitando, según el argumento del film, la caída de la civilización y el final de la edad de bronce. Es una tesis un poco extrema, pero nuestro bienamado director se aferra a ella como a un clavo ardiendo. Quizás por eso suprime de la película uno de los capítulos más importantes de la \emph{Odisea}, ya que Nausícaa y los feacios suponen un contraejemplo a su pesimismo. Y en mi opinión eso es un pecado capital, que puede suponerle varios siglos de purgatorio. No solo elimina uno de los personajes más adorables ---quizás el más adorable, en todos los sentidos--- de la épica de Homero, sino que borra las pruebas que no le convienen a su narrativa un poco estrecha de miras.

Por si no queda claro que los feacios respetan la \gr{ξενία} al pie de la letra, Nausícaa dice:


\begin{verse}
Dijo; y en ese momento llamó a las beltrenzas esclavas:\\
«Alto ahí, muchachas; ¿adónde escapáis porque un hombre nos salga?\\
¿Pensáis acaso que es uno de esos que viene de malas?\\
Aún no respira ese hombre mortal, y tal vez nunca nazca,\\
que de los hombres feacios la tierra venga a tocarla\\
trayendo calamidad: tanto y más los sin muerte nos guardan.\\
Lejos, muy lejos vivimos, en la tempestuosa maralta,\\
el fin del mundo, y con otros mortales no hay trato ni alianzas.\\
Pero nos llega ahora un triste errabundo y debemos sin falta\\
darle atención, pues que toda la gente mendiga o extraña\\
gente es de Zeus, y el regalo que des, aun pequeño, le agrada.\\
Dadle, así, esclavas, a este hombre extranjero bebida y pitanza, y\\
bañadlo bien en el río, y que abrigo del viento allí haya».
\end{verse}

Pero Ulises, pudoroso, le pide a las muchachas que le dejen arreglarse solo. Así lo hace, un poco apartado, y Atenea aprovecha para maquillarlo divinalmente, así que cuando sale a la vista, está hecho un brazo de mar.


\begin{verse}
Fue él a sentarse después, apartándose un poco, a la playa,\\
resplandeciente en gracia y belleza. Ella bien lo miraba.\\
Y, contemplándolo, dijo a las ondulihermosas esclavas:

«Niñas de brazos de leche, oíd lo que voy a deciros.\\
No mal de grado de todos los dioses que alberga el Olimpo\\
viene a encontrarse este hombre aquí con los feacios divinos;\\
que es que creí que era un ser en verdad poco apuesto al principio,\\
y ahora es igual a los dioses, del ancho cielo inquilinos.\\
Ay, si uno así a mí me fuera llamado ya esposo mío\\
viviendo aquí, porque en gusto quedarse por siempre le vino…\\
Ah, dad comida y bebida, esclavas, al desconocido».
\end{verse}

El patito (o más bien el ganso) feo se ha convertido en un apuesto galán por obra y gracia de Atenea, que de paso habrá aprovechado para quitarle alguna cana; tanto, que a la princesa le da por pensar que uno así no sería mal marido.

Pero como tiene tanto de decidida como de hermosa, enseguida informa al náufrago de que se dispone a llevarlo al palacio de su padre y le da instrucciones para entrar en la ciudad. Como no tiene un pelo de tonta, le avisa de que tienen que andarse alerta para evitar las malas lenguas que podrían decir algo como esto:

\begin{verse}
«¿Quién será ese que sigue a Nausícaa, tan grande y garrido\\
y forastero? ¿De dónde salió? Se me da que un marido.\\
Sí, se lo habrá ella llevado cuando él se extravió del navío,\\
que es de otra tierra, y lejana, que aquí no tenemos vecinos;\\
o a ella, tras mucho rezar, algún dios biendeseado le vino\\
caído del cielo, y por todos sus días lo tenga consigo.\\
Tanto mejor si, rodando y rodando, a encontrar fue marido\\
en otra parte, pues que a estos de aquí no los tiene por dignos,\\
a los feacios, y muchos la rondan, y muy distinguidos».
\end{verse}


El tema del matrimonio con el extranjero sigue dando vueltas en el relato, pero también el de la maledicencia, y eso que todavía no se habían inventado las redes sociales:

\begin{verse}
Así dirán, y eso sí para mí que sería un ultraje.\\
Y que lo viera yo en otra, a mí llegaría a indignarme,\\
si ella a pesar de los suyos, viviéndole el padre y la madre,\\
a hombre se ajunta, en lugar de ir a boda, y con pueblo delante.
\end{verse}

Así que la princesa le pide a Odiseo que espere en un bosquecillo cercano para que no los vean llegar juntos. Después ---le explica--- podrá dirigirse al palacio de su padre y pedir audiencia. El canto VI termina de una manera curiosa y muy homérica. No olvidemos que Atenea lleva todo el episodio enredando. Pues bien, cuando Ulises se queda por fin solo en el bosque, le pide que le ayude a hacer amistad con los feacios:

\begin{verse}
«Óyeme, cría de Zeus cabrilámpigo, tú, Intemerenda:\\
préstame ahora el oído que no me prestaste cuando era\\
un hombre roto cuando me rompía el que atrema la tierra.\\
Hazme a los feacios amigo y digno de condolencia».
\end{verse}

¿Y qué hace Atenea? Pues le da por acordarse de que tío Poseidón tiene muy malas pulgas y prefiere, por una vez, estarse quietecita:


\begin{verse}
Dijo rezándole así, y Palas Atena lo oía.\\
Pero a sus ojos no se apareció, pues respeto tenía\\
a ese su tío de padre, tozudo en rencor todavía\\
contra Odiseo pardiós, mientras diera en su tierra algún día.
\end{verse}

En \emph{Abandonando Ítaca}\footnote{\url{https://eolasediciones.es/libro/abandonando-itaca-leaving-ithaca/}}, me pongo en la piel de Odiseo, recordando en Ítaca aquella mañana de encuentros afortunados y aquel momento de soledad en el bosque sagrado de Atenea. A diferencia de Homero, que rara vez se demora en una introspección sostenida, mi viejo Ulises se debate entre la nostalgia y los remordimientos.

\begin{verse}
Cada vez que cierras tus ojos al tenderte en la arena,\\
la ves. A la hermosa princesa,\\
jugando con sus criadas. Todas\\
estaban desnudas, y su belleza te ahogaba,\\
más que el azul desierto del que habías\\
escapado por poco.\\
No salió en estampida como las demás, ni fue tímida. En cambio, tú,\\
que habías violado a muchas vírgenes,\\
apenas podías pronunciar palabra. Es cierto que\\
estabas también desnudo, pero no avergonzado\\
de tu fuerte cuerpo. Tal vez fue algo nuevo,\\
por primera vez viste lo que sus ojos vieron observándote,\\
un mendigo, un náufrago, pidiendo clemencia.
\end{verse}


Y no le queda otra que reconocer unos sentimientos que quizás negó incluso ante sí mismo en aquel momento.


\begin{verse}
¿La amabas? Tal vez sí,\\
pero nunca te atreviste a admitirlo,\\
ni siquiera a ti mismo. En su lugar,\\
noche tras noche seguías desgranando\\
tus largas travesías y aventuras salvajes,\\
el viaje sin fin que te había arrojado\\
junto a sus costas.
\end{verse}


Y cuestionarse ---es el momento de cuestionárselo todo--- si mantuvo su decisión de volver a casa por convicción o cobardía.

\begin{verse}
¿Tomaste en consideración quedarte?\\
Su padre habría estado contento de casar a su hija\\
con el último héroe que quedaba de Troya.\\
Pero decidiste volver,\\
por el bien de Penélope.\\
¿O fue otra la razón?\\
Tal vez sencillamente no pudiste olvidar lo que viste en su rostro:\\
una joven sintiendo compasión por un anciano.
\end{verse}

Nausícaa, como todas las mujeres de la \emph{Odisea}, es polifacética y uno de los ángulos que muestra su historia es el de la posibilidad ---o no--- de una relación amorosa entre una muchacha de quince años y un hombre mayor que le lleva tres o cuatro décadas. Homero se adelanta a Nabokov en tres milenios, pero la princesa de los feacios no es Lolita.

Nausícaa es el amor imposible por partida doble. No solo la pierde el héroe atormentado y ya canoso que arriba a sus costas. La pierde también el adolescente que se quería poner en la piel de Ulises, como si un cordero pudiera suplantar a un león. Nausícaa es tan universal como Helena y Penélope, como Calipso y Circe, y lo es porque todos guardamos en nuestro corazón la imagen de esa muchacha que un día nos dedicó una sonrisa, una mirada, puede que un beso. Todos atesoramos, como un dulce y lacerante filtro, ese amor que pudo ser y no fue. Todos seguimos amando a Nausícaa.
